\subsection{The Fredholm Index}\label{sec:fredholm-index}

For many linear PDEs it is much easier to prove uniqueness than existence. For operators in a finite-dimensional vector space, it is well-known that uniqueness and existence are in fact equivalent; this is known as the Fredholm alternative. It is important to consider those operators in infinite dimensions for which a Fredholm alternative holds. We begin with a definition.

\begin{Definition}
\label{def:semi-fredholm-and-fredholm-operators}
Let $X$ and $Y$ be Banach spaces. We say that the operator $A\in\mathcal{L}(X,Y)$ is \textbf{semi-Fredholm} if $\mathcal{R}(A)$ is closed and if either $\mathcal{N}(A)$ is finite-dimensional or $\mathcal{R}(A)$ has a finite-dimensional complement in $Y$. If both are true, the operator is called \textbf{Fredholm}.
\end{Definition}

We have restricted our definition to bounded operators. However, if $A$ is unbounded, we can always regard it as a bounded operator defined on the Banach space $\mathcal{D}(A)$, where $\mathcal{D}(A)$ is equipped with the graph norm (cf. Problem~\ref{ex:operator-graph-norm-completeness}).

The most important property of semi-Fredholm operators is a quantity called the Fredholm index.

\begin{Definition}
\label{def:fredholm-nullity-deficiency-index}
Let $A\in\mathcal{L}(X,Y)$ be semi-Fredholm. Then the dimension of $\mathcal{N}(A)$ is called the \textbf{nullity} of $A$ and the dimension of the complement to $\mathcal{R}(A)$ is called the \textbf{deficiency} of $A$ (If $\mathcal{R}(A)$ does not have a finite-dimensional complement, the deficiency is infinite.) The quantity
\begin{equation}\label{eq:fredholm-index-nullity-minus-deficiency}
\operatorname{ind}A=\operatorname{nul}A-\operatorname{def}A
\end{equation}
is called the \textbf{(Fredholm) index} of $A$.
\end{Definition}

If $A$ is Fredholm, the index is finite; otherwise it is either plus or minus infinity. The crucial theorem about semi-Fredholm operators is the following:

\begin{Theorem}
\label{thm:semi-fredholm-index-local-stability}
Let $A\in\mathcal{L}(X,Y)$ be semi-Fredholm. Then there exists $\varepsilon>0$ such that any $B\in\mathcal{L}(X,Y)$ with $\|B-A\|<\varepsilon$ is also semi-Fredholm and, moreover, $\operatorname{ind}B=\operatorname{ind}A$.
\end{Theorem}

The proof which we give works if either $A$ is Fredholm or $X$ and $Y$ are Hilbert spaces. The difficulty in the general case is that a closed subspace of a Banach space does not necessarily have a closed complement; we refer to \cite{Ka} for a proof in the general case. For the case when $A$ is Fredholm, we note the following lemma:

\begin{Lemma}
\label{lem:finite-dimension-or-codimension-closed-complement}
Let $X$ be a Banach space and assume that $V$ is a closed subspace of $X$ which is either finite-dimensional or of finite codimension. Then there is a closed subspace $W$ of $X$ such that $X=V\oplus W$.
\end{Lemma}

\begin{Proof}
If $V$ has finite codimension, we merely have to note that every finite-dimensional normed vector space is complete, hence every finite-dimensional subspace of $X$ is closed. If $V$ is finite-dimensional, let $e_i$, $i=1,\ldots,n$, be a basis of $V$. By the Hahn-Banach theorem we can construct linear functionals $x_i^*\in X^*$ such that $x_i^*(e_j)=\delta_{ij}$. We then define $W$ to be the intersection of the nullspaces of the $x_i^*$.
\end{Proof}

In the following proof, we shall also have to use the fact that the direct sum of a closed subspace and a finite-dimensional subspace is closed; we leave the proof of this as an exercise (Problem~\ref{ex:closed-plus-finite-dimensional-subspace}).

We now proceed to the proof of the theorem assuming that either $A$ is Fredholm or $X$ and $Y$ are Hilbert spaces.

\begin{Proof}
Let $V$ and $W$ be subspaces of $X$ and $Y$, respectively, so that $X=\mathcal{N}(A)\oplus V$, $Y=\mathcal{R}(A)\oplus W$. We define $\widetilde A:V\times W\to Y$ by $\widetilde A(v,w)=Av+w$. Then $\widetilde A$ is bijective. Analogously, we define $\widetilde B$ for some given operator $B\in\mathcal{L}(X,Y)$. If $\|B-A\|$ is sufficiently small, then the same is true for $\|\widetilde B-\widetilde A\|$, hence $\widetilde B$ is bijective. In other words, the equation $Bv+w=y$ for given $y\in Y$ has a unique solution $v\in V$, $w\in W$. It follows immediately that $\mathcal{R}(B)+W=Y$, i.e., $\operatorname{def}B\leq\operatorname{def}A$ and that $\mathcal{N}(B)\cap V=\{0\}$, i.e., $\operatorname{nul}B\leq\operatorname{nul}A$. Moreover, $B(V)=\{y\in Y\mid\widetilde B^{-1}(y)\in V\}$ is closed and has finite codimension in $\mathcal{R}(B)$, since either $B(V)$ has finite codimension in $Y$ or $V$ has finite codimension in $X$. Hence $\mathcal{R}(B)$ is closed and $B$ is semi-Fredholm.

Since either $V$ has finite codimension in $X$ or $B(V)$ has finite codimension in $Y$, we conclude that $V+\mathcal{N}(B)$ has finite codimension in $X$. Let now $Z$ be a space such that $X=V\oplus\mathcal{N}(B)\oplus Z$. Then we have $Y=B(V)\oplus W$ and $\mathcal{R}(B)=B(V)\oplus B(Z)$, i.e., $\operatorname{def}B=\dim W-\dim B(Z)=\dim W-\dim Z$ and $\operatorname{nul}B=\dim\mathcal{N}(A)-\dim Z$; the equality of the indices follows immediately.
\end{Proof}

We have the following corollary of Theorem~\ref{thm:semi-fredholm-index-local-stability}:

\begin{Corollary}
\label{cor:semi-fredholm-open-set-constant-index}
The set of all semi-Fredholm operators is open in $\mathcal{L}(X,Y)$. Moreover, the index is constant on each connected component.
\end{Corollary}

Hence, if $A(t)\in\mathcal{L}(X,Y)$ is semi-Fredholm for every $t$ and depends continuously on $t$ for $t\in[a,b]$, then the index is independent of $t$. This is often useful in applications, where the index of $A(a)$ may be easier to find than that of $A(b)$. One application of this approach is the following theorem.
